\section{Introduction}
\label{sec:intro}

\subsection{Motivation: Physics-Native Computing}

Since 1945, the von Neumann architecture has defined digital 
computation: programs are stored in memory and executed by a 
sequential fetch-decode-execute cycle. As Moore's law slows and the 
memory wall tightens, this separation imposes two fundamental limits: 
(a) \emph{energy waste}: moving data costs $100\times$ more than 
computing on it~\cite{horowitz2014computing}; and 
(b) \emph{the clock bottleneck}: complex iterative problems 
(such as Einstein field equations, fluid dynamics, and quantum 
many-body systems) suffer severe latency from sequential iteration 
and ECC verification overhead.

Physics-native computing proposes a different paradigm: rather than 
\emph{simulating} mathematics in hardware, let the hardware 
\emph{be} the mathematics. The intrinsic physical laws of analog 
circuits---Ohm's law, Kirchhoff's laws, capacitor charge dynamics---can 
directly realize variational relaxation and solve constraint-satisfaction 
problems without clocked iteration~\cite{adamatzky2019physics}.

\subsection{Zero-Space Projection: A Brief Recap}

\abnqss v3.0 is a representative physics-native architecture based on 
a silicon-photonics MZI array (L0 layer) coupled with an analog OTA-C 
array (L1 ``Shenqu engine''). Its core theoretical claim is the 
\emph{zero-space projection lock theorem}:
\begin{equation}
    \Cvec_{\mathrm{final}} = \proj \cdot (\Cvec_{\mathrm{base}} 
    - \delta \cdot \alpha \cdot \evec_0),
    \label{eq:projection}
\end{equation}
where $\proj = \Vmat \Vmat^\top$ is the projection matrix onto the 
null space of the constraint matrix $\Mmat$, $\Vmat$ is a null-space 
basis, $\delta$ is a physical perturbation, and $\alpha$ is the 
physical coupling constant.

The digital twin of \abnqss v3.0 has been validated for 
$N = 3, \dots, 6$, showing:
\begin{itemize}[leftmargin=*]
    \item Projection linearity $R^2 = 1.000000$ (exact);
    \item Null-space dimension $K = N(N-1)/2$ (exact);
    \item Physical coupling slope $\alpha \in [0.0075, 0.0115]$, 
          increasing with $N$.
\end{itemize}

\subsection{Motivating the Present Study}

Given the strong theoretical foundation of \ZSP, a natural question 
arises: \emph{where should \ZSP be applied?} 

Multi-object tracking (\MOT) and cooperative perception are natural 
candidate applications, because their core computational step---data 
association---can be cast as a constraint-satisfaction problem. Several 
prior works have advocated for physics-based or analog approaches to 
perception~\cite{...}. 

This paper reports the outcome of a six-month systematic study 
evaluating \ZSP on three perception tasks. \textbf{Our results are 
negative}: in all three tasks, the classical Hungarian algorithm 
is either strictly better (in accuracy, speed, or communication 
cost) or equal. 

\subsection{Contributions}

Our contributions are:
\begin{enumerate}[leftmargin=*]
    \item \textbf{A complete experimental framework} spanning 
          digital-twin verification to real benchmarks, with 
          all code and data released open-source.
    \item \textbf{Systematic negative results} on three perception 
          tasks: single-frame \MOT, multi-frame \MOT, and 
          cooperative perception.
    \item \textbf{A root-cause analysis}: linear assignment problems 
          possess a dual structure that enables exact polynomial 
          algorithms, which \ZSP's iterative nature cannot beat.
    \item \textbf{An explicit applicability boundary}: \ZSP should 
          be applied to problems without good algorithms, not to 
          problems already solved by exact methods.
\end{enumerate}

\subsection{Paper Organization}

\Cref{sec:related} reviews related work. 
\Cref{sec:method} formalizes \ZSP for the three perception tasks. 
\Cref{sec:experiments} presents the experimental results. 
\Cref{sec:analysis} analyzes why \ZSP fails. 
\Cref{sec:discussion} discusses the applicability boundary. 
\Cref{sec:conclusion} concludes.